\paragraph{What all coalition sizes can determine.}
Fix a feature $j$ and let $D\le p$ bound the number of distinct columns in
any path group. Group the increments requiring $j$ by order into the columns
$A_q=\sum_{H\ni j,\,|H|=q+1}a_H$ of a matrix $A$, $0\le q<D$. Write
$v_q(s)=\binom{s}{q}/\binom{p-1}{q}$, $K=A^\top A$ and $g=A^\top\delta_j$.

\begin{proposition}
The directions $h_{js}$ for $s=0,\ldots,D-1$ span the directions for all $p$
sizes. Given the inner products $g$, the smallest compatible squared norm is
\[
 \min_{A^\top x=g}\|x\|^2
   =g^\top K^\dagger g
   \ \ge\ \max_{0\le s<p}Q_{js}.
\]
Equality with the largest individual vote holds exactly when the projection of
$\delta_j$ onto $\operatorname{col}(A)$ is zero or parallel to some nonzero $h_{js}$.
\end{proposition}
\begin{proof}
Equation~\eqref{eq:tree-size} gives $h_{js}=Av(s)$. The first $D$ rows
$v(s)^\top$ form a lower triangular matrix with diagonal
$1/\binom{p-1}{s}>0$. Thus their images span $\operatorname{col}(A)$ even
when the columns of $A$ are dependent. Let $x_*=AK^\dagger g$, where
$K^\dagger$ is the Moore--Penrose inverse. Since $g=A^\top\delta_j$ belongs to
$\operatorname{col}(K)$, we have $A^\top x_*=g$.
Every compatible $x$ differs from $x_*$ by a vector orthogonal to
$\operatorname{col}(A)$. Pythagoras gives the minimum, and Cauchy--Schwarz
on each size direction gives the inequality and its equality condition.
\end{proof}

The bound uses arbitrary signed mixtures of size directions. It coincides with
an individual vote when the span is one-dimensional. For a strict example
take $p=3$, $a_{\{1\}}=(1,0)$, $a_{\{1,2\}}=(0,-1)$ and
$a_{\{1,3\}}=(0,2)$, with all other increments zero. After removing feature 3,
$\delta_1=(1,-1)$ and $h_{1s}=(1,s/2)$. The largest of the three votes is one
and the span bound is two. Sharpness here concerns the stated inner products.
To condition also on the cross term in~\eqref{eq:tree-error}, append $r$ to
$A$ and $(c_j-\|r\|^2)/2$ to $g$.

Individual size votes can also be checked from $g$ and $K$. For remaining
allowance $t=b-c_j\ge0$, a veto occurs exactly when some integer $0\le s<p$ has
$(v(s)^\top g)^2-t\,v(s)^\top Kv(s)>0$. Each $v_q$ is a degree-$q$
polynomial in $s$, so the tested polynomial has degree at most $2D-2$.
A zero direction contributes zero. If $t<0$, the known cross term already
exceeds the cutoff.
